\noindent\textbf{Einheit 1 · Prozente als Anteile}

\erg{1}{a) $20\,\%$ \quad b) $60\,\%$ \quad c) $10\,\%$ \quad d) $90\,\%$ \quad e) 4 der 10 Kästchen grau \quad f) 6 Kästchen und ein halbes grau}
\erg{2}{a) $23\,\%$ \quad b) $9\,\%$ \quad c) $57\,\%$ \quad d) $80\,\%$ \quad e) $52\,\%$ \quad f) $36\,\%$ \quad g) $0{,}62$ \quad h) $\frac{2}{5}$ \quad i) $10\,\%$ \quad j) „16 von 100 Schrauben sind zu kurz.“}
\erg{3}{a) $8\,\%$, $18\,\%$, $80\,\%$ \quad b) $\frac{1}{5}$, $25\,\%$, $0{,}3$ \quad c) Packung B, denn $\frac{2}{5} = 40\,\%$}
\erg{4}{a) $60\,\%$ \quad b) $75\,\%$ \quad c) $30\,\%$ \quad d) $25\,\%$ \quad e) Limonade $35\,\%$ \quad f) Limonade}
\erg{5}{a) Jonas hat den Nenner als Prozentzahl gelesen; jeder achte ist $\frac{1}{8} = 0{,}125 = 12{,}5\,\%$ \quad b) $5\,\%$}
\erg{6}{a) $20 \cdot 5 = 100$, also $\frac{1}{20} = \frac{5}{100} = 5\,\%$ \quad b) nein: $0{,}5\,\% = \frac{0{,}5}{100} = 0{,}005$, dagegen ist $0{,}5 = 50\,\%$}
\erg{7}{a) der zweite Satz (Bus) \quad b) z.\,B. „Drei Zehntel der Schüler kommen mit dem Bus.“ – $30\,\%$; ein Drittel wären rund $33\,\%$}
